\section{Summary and outlook}
\label{sec:summary}

\paragraph{Summary.}
Unbinned OmniFold with gradient-boosted trees reproduces the published
MINERvA ME FHC CC-inclusive $d^{2}\sigma/(d\pT\,d\pPar)$
and its uncertainty scale (\S\ref{sec:execsummary}, \S\ref{sec:results}). The
result passes closure, iteration, pull, calibration, and classifier-family
checks, but the existing 2D toys do not furnish an independent-truth coverage
test. On that foundation we
performed \emph{unbinned} simultaneous three-, four- and five-observable
unfolds of MINERvA data ($+\Eavail$, $+q_{3}$, $+W$), each validated against its
lower-dimensional predecessor by projection anchors agreeing to $\approx$1--2\,\%. (Binned MINERvA multi-differential
measurements, up to triple-differential~\cite{MINERvA:2022qe}, exist; the methodological distinction is the
unbinned, simultaneously-unfolded formulation.) The physics output:
\begin{itemize}[itemsep=2pt]
  \item the low-$\Eavail$ excess over all four truth-level generators
        recovers the established MINERvA low-recoil/2p2h deficit through an
        independent method, with the expected FSI insensitivity and Valencia
        2p2h response (\S\ref{sec:3d-2p2h});
  \item the $(\Eavail,q_{3})$ plane resolves the deficit's momentum-transfer
        structure (\S\ref{sec:4d}); (the 4D systematic-budget composition is not
        quoted --- no corrected 4D covariance is adopted, \S\ref{sec:3d});
  \item the unified-throw machinery now tests the block sum with asymmetric
        endpoints, a fixed estimator seed, and mean-centered joint throws;
        the resulting 5D covariance remains a candidate until the
        selection-complete lateral replacement lands; no corrected 4D covariance is reported
        (\S\ref{sec:uq-combine});
  \item the $(\Eavail,W)$ study localizes the remaining high-$\Eavail$
        excess to the deep-inelastic region, where all four generators
        underpredict the central result; no significance is quoted without a
        corrected projected covariance
        (\S\ref{sec:eavailw});
  \item the recoil-only PET point-cloud track provides a low-level-input
        representation cross-check, yielding an absolutely-normalized central
        value that passes ordinary closure. Its historical five-component budget
        is quarantined because nuisance and retraining shifts require a joint
        covariance (or explicit cross term) and the detector sample was
        CV-support-limited (\S\ref{sec:pet}). Because the muon is omitted from both classifiers,
        these products do not establish a joint muon--hadron unfold and no
        covariance component transfers automatically to the full-event replacement.
        The whole track is diagnostic and method development: no PET covariance
        is adopted as a publication uncertainty product, and the full-event
        central/statistical pairing was declined rather than resolved.
\end{itemize}

\paragraph{Outlook.}
The scalar extended-fiducial study (\S\ref{sec:fps}) yields a preliminary
central value $\sigma_{\mathrm{ext}}=\SI{\fpsSigma}{cm^2/nucleon}$
(\SI{\fpsAbove}{\percent} above the restricted phase space), with an anchor,
closure tests and a three-prior envelope but without a corrected covariance
(\S\ref{sec:uq-combine}). About \SI{6}{\percent} of the rate is a newly
acceptance-supported extension, while about \SI{28}{\percent} is a separately
labeled prior-dominated extrapolation. A full-event PET estimator would add the
distinguished muon and validated event context, require an omitted-muon stress
closure, and repeat the nominal fit and complete UQ on the canonical extended
grid; that is method development and is not on the publication critical path,
which runs through the adopted scalar five-dimensional covariance. Natural
additional refinements are NEUT as a fifth generator and the finer
44-cell Ascencio comparison (the maximal-common-grid
comparison is presented in \S\ref{sec:4d}). One
paper-independent cross-check would further strengthen the analysis --- a
flux\,$\leftrightarrow$\,muon
posterior from MAT-MINERvA for a fully self-derived $\rho$. The
truncated-spectral ``ours-only'' goodness-of-fit convention
(App.~\ref{sec:rank}) is no longer among them: the collaboration confirmed it
uses the same truncated-spectral pseudo-inverse.
